% Generated draft: requires mathematical and source review.
\documentclass[11pt]{article}
\usepackage[margin=25mm]{geometry}
\usepackage{fontspec}
\setmainfont{Latin Modern Roman}
\usepackage{amsmath,amssymb,mathtools,unicode-math}
\setmathfont{Latin Modern Math}
\usepackage{enumitem,needspace,xurl,hyperref}
\setlength{\parindent}{0pt}
\setlength{\parskip}{5pt}
\setlength{\emergencystretch}{3em}
\newcommand{\N}{\mathbb{N}}
\newcommand{\Z}{\mathbb{Z}}
\newcommand{\Q}{\mathbb{Q}}
\newcommand{\R}{\mathbb{R}}
\newcommand{\C}{\mathbb{C}}
\newcommand{\F}{\mathbb{F}}
\newcommand{\A}{\mathbb{A}}
\newcommand{\PP}{\mathbb{P}}
\newcommand{\E}{\mathbb{E}}
\newcommand{\eps}{\varepsilon}
\newcommand{\dd}{\,\mathrm{d}}
\begin{document}
\section*{UnsolvedMath ID: 800130}
\textbf{Source problem code:} catalogue entry
\subsection{Definitions and mathematical statement}
Let \(X=B\Gamma_H\) denote the space appearing in the supplied source. The record does not define \(\Gamma_H\) or specify the construction of \(B\Gamma_H\); consequently, identifying the intended space remains an unresolved prerequisite for a standalone formulation. Write \(H_2(X;\Z)\) for its second singular homology group with integer coefficients. For a closed oriented surface \(\Sigma\) and a continuous map \(f\colon\Sigma\to X\), the associated surface cycle is \(f_*[\Sigma]\in H_2(X;\Z)\), where \([\Sigma]\) is the fundamental class.

The source supplies the background assertions
\[
\pi_2(X)=0,
\qquad
f_*[T^2]=0\quad\text{for every continuous }f\colon T^2\to X,
\]
where \(T^2=S^1\times S^1\). It also states that there are no nonzero spherical cycles, meaning that \(g_*[S^2]=0\) for every continuous map \(g\colon S^2\to X\). These assertions are retained as supplied background, without independent verification.

The question, attributed in the source to Tsuboi, is:
\[
\text{Is }H_2(B\Gamma_H;\Z)=0\text{?}
\]
The intended meaning of \(B\Gamma_H\) must be recovered from the cited source before this question can be regarded as fully specified.
\subsection{Short English statement}
Determine whether the second integral homology of the source's space \(B\Gamma_H\) vanishes. The source states that spherical and torus cycles vanish, but the supplied record does not define the space itself.
\subsection{Sources}
\begin{itemize}
\item Local UnsolvedMath registry: permanent ID \(800130\), exact source problem code \texttt{catalogue entry}, title \emph{Foliations and Diffeomorphism Groups --- Question 71}. The supplied registry labels the problem open; this label is not an independently verified mathematical or literature assessment.
\item Nariman, Yazdi, \emph{Problem List on Foliations and Diffeomorphism Groups} (2025), Question 71. Supplied source URL: \url{https://arxiv.org/abs/2503.18274}. The record reports access on 2026-10-16 and extraction by \texttt{arxiv\_latex}; these are supplied provenance details, not a claim of independent access.
\item The question is attributed within the supplied statement to Tsuboi, with reference identifier \texttt{MR3726712}. No full bibliographic entry for that reference is supplied.
\end{itemize}
\end{document}
